% Imporditud: tools/import_eksperimendid.py
% Allikas: Eesti füüsikaolümpiaadi eksperimendi ülesannete
%   kogu (Rael Kalda, 2023), ülesanne Ü31, lahendus L31.
\setAuthor{Konstantin Dukats}
\setRound{piirkonnavoor}
\setYear{2022}
\setNumber{P E2}
\setDifficulty{1}
\setTopic{vedelike mehaanika}
\setVahendid{klaas veega, joonlaud, plastiliin}
\setPunktid{10}
\setAllikas{Eksperimendi kogu Ü31, lahendus lk 46}
\setLahendusseis{toores}

\prob{Plastiliin}
Leidke plastiliini tihedus. Vee tihedus on $\rho$ = 1000 kg m$^{-3}$.

\hint

\solu
Määrata plastiliini ruumala sukeldumismeetodil. Valmistada plastiliinist laevuke ja mõõta selle poolt väljatõrjutud vedeliku ruumala (3 p). Tuletada ujumise tingimusest valem plastiliini tiheduse arvutamiseks:

$\rho$p = mp = $\rho$v Vv (3 p) Vp Vp

Tuleb teha vajalikud mõõtmised ja arvutada plastiliini tihedus (3 p). Vastus oleneb plastiliini sordist ja on ca 1,1...1,2 g cm$-$3 (1 p).
\probend
